\section{Densities and planar tangent measures}\label{sec:density}

We isolate the geometric statement that will be used to finish
the proof. In its hypotheses, \(P\) need not annihilate closed
invariant forms.

\begin{theorem}[Closedness of the positive-density part]
\label{thm:splitting}
Fix a smooth almost complex structure \(J\) and a \(J\)-Hermitian
metric on a closed four-manifold. Suppose that
\[
P(\alpha)=\int_{\mathscr C}\alpha_x(\ell)\dd\lambda(x,\ell)
\]
for a finite positive measure \(\lambda\), with projection \(\mu\).
Suppose also that \(Q\in L^2\) is anti-invariant and that \(P+Q\)
is a closed current. Assume the estimates
\begin{align}
\mu(B_r(x))&\le Cr^2,\label{eq:abstract-growth}\\
\iint_{d_g(x,y)<r}|\ell-\tau_{yx}m|^2
\dd\lambda(x,\ell)\dd\lambda(y,m)&\le Cr^4
\label{eq:abstract-angular}
\end{align}
for all sufficiently small \(r\), uniformly in \(x\).
Then the limit
\[
\theta(x)=\lim_{r\downarrow0}r^{-2}\mu(B_r(x))
\]
exists at every point. If \(E=\{x:\theta(x)>0\}\), the
current \(P_E=\ind_EP\) is closed.
\end{theorem}

For positive currents that are already closed, Elkhadhra proves
that restriction to a \(J\)-analytic subset remains closed
\cite[Section~3.2, Lemma~1, author version]{Elkhadhra2014}.
Here only \(P+Q\) is initially closed, and the positive-density
set is not assumed to be \(J\)-analytic; the mass and angular
estimates provide the cutoff control needed for the restriction.

Throughout this section and Section~\ref{sec:splitting}, assume
the hypotheses of Theorem~\ref{thm:splitting}, and write \(T=P+Q\).
We continue to write \(\ang_r\) for the left side of
\eqref{eq:abstract-angular}. We first establish the density
limit and show that, at \(\mu\)-almost every positive-density point,
the current has planar tangent measures. This makes the average
transverse displacement small without supplying a rate.
Section~\ref{sec:splitting} combines that smallness with the
angular bound to prove closedness.

\begin{lemma}\label{lem:density}
The two-density \(\theta(p)\) exists and is finite at every
\(p\in X\).
\end{lemma}

The radial comparison is in the tradition of Lelong--Jensen
formulas for almost complex currents; see
\cite[Proposition~1 and Corollary~2]{ElkhadhraMimouni2007}.
We apply closure to \(P+Q\) and estimate the \(L^2\) correction
directly.

\begin{proof}
The current \(T\) has measure coefficients with variation
bounded by a constant times
\(\sigma=\mu+|Q|\,dV_g\). Cauchy--Schwarz and
\eqref{eq:abstract-growth} give
\[
\sigma(B_s(p))\le Cs^2,
\]
uniformly in \(p\).

In normal coordinates \(z\) at \(p\), let \(J_0=J(p)\) and
\(F_0=F_p\) be constant, and put \(t=|z|=d_g(p,y)\).
Our sign convention gives
\(dd^c_{J_0}(t^2/2)=2F_0\). Call a radius \emph{good} if its sphere is
\(\sigma\)-null. For good radii define
\[
b(r)=r^{-2}T(\ind_{B_r(p)}F_0).
\]
These pairings are well-defined using the coefficient
measures. We claim that
\begin{equation}\label{eq:density-identity}
b(s)-b(r)=\tfrac12
T\bigl(\ind_{\{r<t<s\}}dd^c_{J_0}\log t\bigr),
\qquad 0<r<s.
\end{equation}

To verify it, define
\[
\phi_a(t)=
\begin{cases}
\log a+(t^2-a^2)/(2a^2),&t<a,\\
\log t,&t\ge a.
\end{cases}
\]
The function \(\phi_r-\phi_s\) is compactly supported and
\(C^{1,1}\). Its \(dd^c_{J_0}\) is
\(2(r^{-2}-s^{-2})F_0\) on \(B_r\), equals
\(dd^c_{J_0}\log t-2s^{-2}F_0\) on the annulus, and is
zero outside \(B_s\). For fixed \(r,s\), compactly supported
smooth approximations have uniformly bounded Hessians and
converge in second derivatives off the two spheres.
Dominated convergence against \(\sigma\) permits testing
closure on their exact two-forms and passing to the limit.
The resulting identity is \eqref{eq:density-identity}.

The form \(dd^c_{J_0}\log t\) is \(J_0\)-invariant and
semipositive, with size \(O(t^{-2})\).
Since \(J(y)-J_0=O(t)\), it evaluates on actual complex
line bivectors to at least \(-C/t\), and its
\(J\)-anti-invariant projection has size \(O(1/t)\).
The latter estimate controls its pairing with \(Q\).
Inward dyadic shells and the quadratic bound on \(\sigma\)
give
\[
\int_{\{0<t<s\}}t^{-1}\dd\sigma\le Cs.
\]
It follows that
\begin{equation}\label{eq:almost-monotone}
b(s)-b(r)\ge-Cs.
\end{equation}
Also \(T(\ind_{B_r}F)=\mu(B_r)\), since
\(\ip{Q}{F}_g=0\) almost everywhere. The bound
\(F-F_0=O(t)\) therefore implies
\begin{equation}\label{eq:density-error}
|b(r)-r^{-2}\mu(B_r(p))|\le Cr.
\end{equation}

In particular \(b\) is bounded. Fix a good radius \(s\)
in \eqref{eq:almost-monotone} and let \(r\downarrow0\)
through good radii; then let \(s\downarrow0\) along a
sequence realizing the lower limit of \(b\).
The upper and lower limits agree. Equation
\eqref{eq:density-error} gives the density limit along
good radii.

Only countably many spheres at a fixed center can have
positive \(\sigma\)-mass. Any radius \(r\) can therefore
be bracketed by good radii in
\(((1-\eps)r,r)\) and \((r,(1+\eps)r)\).
Ball inclusion, followed by \(r\downarrow0\) and then
\(\eps\downarrow0\), proves the asserted limit along
all radii. Its finiteness follows from
\eqref{eq:abstract-growth}.
\end{proof}

For \(d_g(x,y)\) below the injectivity radius, set
\(z_x(y)=\exp_x^{-1}(y)\).

\begin{lemma}[Vanishing transverse moment]\label{lem:transverse}
Under the hypotheses of Theorem~\ref{thm:splitting},
\begin{equation}\label{eq:transverse}
\trans_r:=
\int_{\mathscr C}\int_{B_r(x)}
\left|\pr_{L_\ell^\perp}\frac{z_x(y)}r\right|^2
\dd\mu(y)\dd\lambda(x,\ell)=o(r^2).
\end{equation}
\end{lemma}

\begin{proof}
Write \(P=A\mu\) in a smooth local frame, so that \(A(p)\)
is the conditional average of the line bivectors over \(p\).
At \(\mu\)-almost every point, differentiation with respect
to the Radon measure \(\mu\) gives
\begin{equation}\label{eq:polarization-differentiation}
r^{-2}\int_{B_{Dr}(p)}|A(y)-A(p)|\dd\mu(y)\longrightarrow0
\qquad(D<\infty\text{ fixed}).
\end{equation}
For the precise differentiation input one may use
\cite[Theorem~4.33 and Corollary~4.34]{Kinnunen2026}.
To see the normalization, enclose \(B_{Dr}(p)\) in a
Euclidean coordinate ball of radius \(CDr\).
The mean oscillation on that ball tends to zero, while
its mass divided by \(r^2\) is bounded by
\eqref{eq:abstract-growth}. This does not require
\(\mu\) to be doubling. Smooth changes of frame add a
vanishing \(O(r)\) error after this normalization.

Fix such a differentiation point \(p\) with \(\theta(p)>0\).
The locally rescaled measures
\[
\nu_r=r^{-2}(z_p/r)_*\mu
\]
have uniformly bounded mass on each compact subset of
\(T_pX\), so every sequence of scales has a vaguely
convergent subsequence. Every limit \(\nu\) satisfies
\begin{equation}\label{eq:tangent-ball}
\nu(B_D(0))=\theta(p)D^2\qquad(D>0).
\end{equation}
Initially the identity follows at continuity radii of
\(\nu\). Increasing continuity radii recover the mass of
each open ball, while decreasing continuity radii recover
the corresponding closed ball. Both masses are
\(\theta(p)D^2\), so the identity holds for every \(D>0\)
and every centered sphere is \(\nu\)-null.

The current and measure rescalings have the same normalization.
Indeed, put \(D_r(y)=z_p(y)/r\). For a compactly supported
test two-form \(\beta=\sum_{i<j}\beta_{ij}(w)\,dw_i\wedge dw_j\)
on \(T_pX\),
\[
D_r^*\beta=r^{-2}\sum_{i<j}\beta_{ij}(z_p/r)\,dz_i\wedge dz_j.
\]
Thus the pushforward \((D_r)_*T\), defined by testing \(T\)
against \(D_r^*\beta\), already has the \(r^{-2}\) factor
in its coefficient measures; no further normalization is used.
Equation
\eqref{eq:polarization-differentiation} makes its \(P\)
part converge to \(A(p)\nu\). Its \(Q\) part vanishes
in variation on each fixed ball, since
\begin{equation}\label{eq:q-blowup}
r^{-2}\int_{B_{Dr}(p)}|Q|\vol
\le C_D\|Q\|_{L^2(B_{Dr}(p))}\longrightarrow0.
\end{equation}
The last limit holds at every point by absolute continuity
of the integral of \(|Q|^2\). Compactly supported tests
in the rescaled coordinates pull back inside the original
coordinate chart for small \(r\). Thus closure passes to
the limit, and \(A(p)\nu\) is a closed constant-polarization
current on \(T_pX\).

For a constant skew matrix \(A\), this closure says
\(\sum_i A^{ij}\partial_i\nu=0\) for each \(j\).
Consequently \(\partial_v\nu=0\) for every \(v\in\im A(p)\).
For a compactly supported smooth function \(\phi\), the derivative
of \(t\mapsto\int\phi(w+tv)\dd\nu(w)\) is therefore zero.
This proves that \(\nu\) is invariant under translations along
\(\im A(p)\). The normalization \(F_p(A(p))=1\) excludes
rank zero. Rank four would imply that \(\nu\) is a
constant multiple of four-dimensional Lebesgue measure,
contrary to \eqref{eq:tangent-ball}.
Therefore \(A(p)\) has real rank two.

The associated positive Hermitian matrix has trace one
and complex rank one. Thus \(A(p)\) is itself a unit
complex line bivector. If \(\lambda_p\) is the conditional
probability measure on \(\mathscr C_p\), then
\[
\int_{\mathscr C_p}|\ell-A(p)|^2\dd\lambda_p(\ell)
=1-|A(p)|^2=0.
\]
Hence \(\lambda_p\) is concentrated on \(A(p)\).
Denote its plane by \(L\).

Translation invariance along \(L\) implies
\[
\nu=\cL^2_L\otimes\gamma
\]
for a positive Radon measure \(\gamma\) on \(L^\perp\).
For example, testing against a nonnegative compactly
supported function in the normal variables gives a
translation-invariant measure in the tangential variables,
hence a multiple of area; these multiples define \(\gamma\).
Dividing the centered ball mass by \(D^2\) in
\eqref{eq:tangent-ball} yields
\begin{equation}\label{eq:transverse-product}
\theta(p)=\pi\int_{L^\perp}(1-|w|^2/D^2)_+\dd\gamma(w).
\end{equation}
For \(D_2>D_1\), the difference of these integrands is
nonnegative and strictly positive at every
\(0<|w|<D_2\). Since the integrals are equal, \(\gamma\)
is supported at the origin. Every tangent limit is
therefore \((\theta(p)/\pi)\cL^2_L\).

It follows that the normalized inner integral in
\eqref{eq:transverse} tends to zero at
\(\lambda\)-almost every \((p,\ell)\) of positive density.
Indeed, in the rescaled coordinates a continuous cutoff
equal to one on \(B_1\), supported in \(B_2\), times
\(|\pr_{L^\perp}z|^2\) dominates the integrand and
vanishes on every tangent limit. Subsequential compactness
gives convergence along all scales.

When \(\theta(p)=0\), the same conclusion follows directly
by bounding the normalized inner integral by
\(r^{-2}\mu(B_r(p))\). These normalized integrals are
uniformly bounded by \eqref{eq:abstract-growth}.
Dominated convergence against the finite measure
\(\lambda\) proves \eqref{eq:transverse}.
\end{proof}
